\documentclass[10pt,a4paper]{amsart}
\usepackage[margin=19mm]{geometry}
\usepackage{amsmath,amssymb,mathtools}
\usepackage[scr=rsfs,cal=euler]{mathalfa}
\usepackage{xfrac}
\usepackage[unicode,hidelinks,bookmarks=false]{hyperref}
\newcommand{\ie}{i.e.\ }
\newcommand{\cf}{cf.\ }
\newcommand{\resp}{respectively\ }
\newcommand{\Subsection}{\subsection}
\providecommand{\nobreakdash}{\nobreak\hbox{-}}
\setlength{\emergencystretch}{2em}
\multlinegap0pt
\usepackage{amssymb}
\usepackage{enumerate}
\usepackage{tikz}
\usepackage{tcolorbox}
\newtheorem{proposition}{Proposition}
\newtheorem{lemma}{Lemma}
\title{Uniqueness of Ground State Solutions for a Defocusing Hartree Equation via Inverse Optimal Problems}
\author{Yavdat~Il'yasov$^{a}$, Juntao Sun$^{b}$ , Nur~Valeev$^{a}$, Shuai Yao$^{b}$}
\date{Source: arXiv:2601.14591v1 (CC BY 4.0)}
\begin{document}
\maketitle

\noindent\emph{Source and licence.} Uniqueness of Ground State Solutions for a Defocusing Hartree Equation via Inverse Optimal Problems. Version arXiv:2601.14591v1 (CC BY 4.0), distributed by arXiv under the Creative Commons Attribution 4.0 International licence, http://creativecommons.org/licenses/by/4.0/, as recorded in the arXiv licence metadata for this exact version and shown on the abstract page https://arxiv.org/abs/2601.14591v1. Full licence notice: \texttt{licenses/arxiv-2601.14591-CC-BY-4.0.txt}.\par\smallskip

\noindent\textit{Editorial scope.} Counted unit: the article's lemma lem51 (source0.tex lines 1208--1218). The article's complete original proof of it is delivered with it (source0.tex lines 1220--1323, one \texttt{proof} environment of 104 lines). No mathematics is written, repaired or re-derived: the statement and the proof are the article's own lines, in the article's own notation, and no retained line is substituted or re-flowed.  Retained uncounted as the source-local context the proof needs: lines 253--258; lines 601--604; lines 809--814.  Nothing was omitted from any retained span, so the package carries no omission record.  No retained line contains a citation, so the wrapper carries no bibliography and no bibliography entry.  No retained line contains a figure environment, an image-inclusion command or drawing code, so no image is delivered and the package declares no asset dependency.  Editorial formatting only, in addition: an AMS \texttt{amsart} preamble that restates the article's own declarations and macros used by the retained lines, namely the environments \texttt{proposition}, \texttt{lemma} on separate counters, together with the standard packages those lines need.  The retained environments print in this document as Proposition 1, Lemma 1 in order of appearance, whereas the article prints the same items with the numbering of its own full document.  The complete native source of the article as distributed in the arXiv e-print for this exact version is bundled unchanged as \texttt{sources/193099/source0.tex}.
\begin{equation}
\hat{\mathcal{P}}(\lambda ,\bar{\rho}):=\Vert \bar{\rho}-\hat{\rho}\Vert _{%
\mathbb{L}_{\mu }^{2}}^{2}=\min \left\{ \Vert \bar{\rho}-\rho \Vert _{%
\mathbb{L}_{\mu }^{2}}^{2}\mid \rho \in \mathbb{L}_{\mu }^{2},\ \lambda
=\lambda _{1}(\rho )\right\} .  \label{eq:p}
\end{equation}%

\begin{proposition}
\label{prop2} $\lambda_1(\rho)$ is a upper semicontinuous functional in $%
\mathbb{L}_{\mu}^{2}$.
\end{proposition}

\begin{lemma}
\label{lem1} For any $\bar{\rho}\in \mathbb{L}_{\mu }^{2}$ and $\lambda
>\lambda _{1}(\bar{\rho})$ the minimization problem (\ref{eq:p}) admits a
unique solution $\hat{\rho}(\bar{\rho})\in \mathbb{L}_{\mu }^{2}\setminus
\{0\}$.
\end{lemma}

\begin{lemma}
\label{lem51}

(i) The map $\mathbb{L}^{2}_{\mu} \ni \bar{\rho} \mapsto \hat{\rho}(\lambda,%
\bar{\rho}) \in \mathbb{L}^{2}_{\mu}$ is continuous for any $\lambda \geq
\lambda_1(\bar{\rho})$;

(ii) The map $[\lambda _{1}(\bar{\rho}),+\infty )\ni \lambda \mapsto \hat{%
\rho}(\lambda ,\bar{\rho})\in \mathbb{L}_{\mu }^{2}$ is continuous for any $%
\bar{\rho}\in \mathbb{L}_{\mu }^{2}$.
\end{lemma}

\begin{proof}
Let us prove (i). Assume that $\lambda \in (\lambda _{1}(\bar{\rho}),+\infty
)$, and $\rho _{n}\rightarrow \bar{\rho}$ in $\mathbb{L}_{\mu }^{2}$ as $%
n\rightarrow \infty $. Recall that, by Proposition~\ref{prop2}, the
functional $\lambda _{1}(\cdot )$ is upper semicontinuous on $\mathbb{L}%
_{\mu }^{2}$. Hence $\lambda >\lambda _{1}(\rho _{n})$ for all sufficiently
large $n$, and therefore, by Lemma \ref{lem1} there exists a unique
minimizer $\hat{\rho}(\rho _{n})$ of
\begin{equation*}
\hat{\mathcal{P}}(\lambda ,\rho _{n}):=\Vert \rho _{n}-\hat{\rho}(\rho
_{n})\Vert _{\mathbb{L}_{\mu }^{2}}^{2}=\min \left\{ \Vert \rho _{n}-\rho
\Vert _{\mathbb{L}_{\mu }^{2}}^{2}\text{ }|\text{\ }\rho \in \mathbb{L}_{\mu
}^{2},\text{ }\lambda _{1}(\rho )\geq \lambda \right\} .
\end{equation*}

Using that $\lambda _{1}(\rho )$ is upper semicontinuous on $\mathbb{L}_{\mu
}^{2}$, it follows similarly to the proof of Proposition \ref{prop2} that $%
\hat{\mathcal{P}}(\lambda ,\cdot )$ is an upper semicontinuous functional on
$\mathbb{L}_{\mu }^{2}$ and therefore
\begin{equation}
\hat{\mathcal{P}}(\lambda ,\bar{\rho})\geq \limsup_{n\rightarrow +\infty }%
\hat{\mathcal{P}}(\lambda ,\rho _{n})\equiv \limsup_{n\rightarrow +\infty
}\Vert \hat{\rho}(\rho _{n})-\rho _{n}\Vert _{\mathbb{L}_{\mu }^{2}}.
\label{eq:mathcaP_conv}
\end{equation}%
Hence, the sequence $(\hat{\rho}(\rho _{n}))$ is bounded in $\mathbb{L}_{\mu
}^{2}$, and therefore, the reflexivity of $\mathbb{L}_{\mu }^{2}$ yields
that there exists a subsequence, still denoted by $(\hat{\rho}(\rho _{n}))$,
such that $\hat{\rho}(\rho _{n})\rightharpoonup \check{\rho}$ weakly in $%
\mathbb{L}_{\mu }^{2}$ for some $\check{\rho}\in \mathbb{L}_{\mu }^{2}$.
Analyzing similarly to the proof of Lemma \ref{lem1} we conclude that $%
\check{\rho}\neq 0$ and $\check{\rho}\in \tilde{M}_{\lambda }$.

It is easily seen that
\begin{equation}
\Vert \check{\rho}-\bar{\rho}\Vert _{\mathbb{L}_{\mu }^{2}}\leq
\liminf_{n\rightarrow +\infty }\Vert \hat{\rho}(\rho _{n})-\rho _{n}\Vert _{%
\mathbb{L}_{\mu }^{2}}.  \label{eq:52}
\end{equation}%
Hence by (\ref{eq:mathcaP_conv}) we derive
\begin{equation*}
\Vert \check{\rho}-\bar{\rho}\Vert _{\mathbb{L}_{\mu }^{2}}\leq \hat{%
\mathcal{P}}(\lambda ,\bar{\rho}).
\end{equation*}%
Since $\check{\rho}\in \tilde{M}_{\lambda }$, equality must hold: $\Vert
\check{\rho}-\bar{\rho}\Vert _{\mathbb{L}_{\mu }^{2}}=\hat{\mathcal{P}}%
(\lambda ,\bar{\rho})$. By the uniqueness of the minimizer of (\ref{eq:p}),
this implies that $\check{\rho}=\hat{\rho}(\bar{\rho})$. Moreover, (\ref%
{eq:mathcaP_conv}) and (\ref{eq:52}) yield
\begin{equation*}
\Vert \hat{\rho}(\bar{\rho})-\bar{\rho}\Vert _{\mathbb{L}_{\mu
}^{2}}=\lim_{n\rightarrow +\infty }\Vert \hat{\rho}(\rho _{n})-\rho
_{n}\Vert _{\mathbb{L}_{\mu }^{2}},
\end{equation*}%
and thus,
\begin{equation*}
\hat{\rho}(\rho _{n})\rightarrow \hat{\rho}(\bar{\rho})\quad \text{in}\quad
\mathbb{L}_{\mu }^{2}.  \label{eq:53}
\end{equation*}%
Since $\hat{\rho}(\bar{\rho})$ is uniquely defined, this implies that the
map $\mathbb{L}_{\mu }^{2}\ni \bar{\rho}\mapsto \hat{\rho}(\lambda ,\bar{\rho%
})\in \mathbb{L}_{\mu }^{2}$ is continuous.

Now we prove (ii). Assume that $\bar{\rho}\in \mathbb{L}_{\mu }^{2}$ and $%
\lambda >\lambda _{1}(\bar{\rho})$, and $\lambda _{n}\rightarrow \lambda $
as $n\rightarrow \infty $. Then $\lambda _{n}>\lambda _{1}(\bar{\rho})$ for
all sufficiently large $n$. Hence, by Lemma \ref{lem1} there exists a unique
minimizer $\hat{\rho}(\lambda _{n}):=\hat{\rho}(\lambda _{n},\bar{\rho})$ of
\begin{equation*}
\hat{\mathcal{P}}(\lambda _{n},\bar{\rho}):=\Vert \bar{\rho}-\hat{\rho}%
(\lambda _{n})\Vert _{\mathbb{L}_{\mu }^{2}}^{2}=\min \left\{ \Vert \bar{\rho%
}-\rho \Vert _{\mathbb{L}_{\mu }^{2}}^{2}\text{ }|\text{\ }\rho \in \mathbb{L%
}_{\mu }^{2},\text{ }\lambda _{1}(\rho )\geq \lambda _{n}\right\} .
\end{equation*}%
As before, it can be shown that there exists a weak limit point $\breve{\rho}%
\in \mathbb{L}_{\mu }^{2}\setminus \{0,\bar{\rho}\}$ of the sequence $(\hat{%
\rho}(\lambda _{n}))$ in $\mathbb{L}_{\mu }^{2}$, and
\begin{equation*}
\hat{\mathcal{P}}(\lambda ,\bar{\rho})\geq \limsup_{n\rightarrow +\infty }%
\hat{\mathcal{P}}(\lambda _{n},\bar{\rho}).
\end{equation*}%
Therefore,
\begin{equation*}
\Vert \breve{\rho}-\bar{\rho}\Vert _{\mathbb{L}_{\mu }^{2}}\leq
\liminf_{n\rightarrow +\infty }\Vert \hat{\rho}(\lambda _{n})-\bar{\rho}%
\Vert _{\mathbb{L}_{\mu }^{2}}\leq \hat{\mathcal{P}}(\lambda ,\bar{\rho}).
\label{eq:58}
\end{equation*}

Since the functional $\lambda(\cdot)$ is upper semicontinuous on $\mathbb{L}%
_{\mu}^{2}$, it follows that
\begin{equation*}
\lambda(\breve{\rho})\geq \limsup_{n\to+\infty}\lambda(\hat{\rho}%
(\rho_n))\geq \lambda.
\end{equation*}
Thus, $\breve{\rho} \in \tilde{M}_{\lambda}$, and therefore
\begin{equation*}
\|\breve{\rho} -\bar{\rho}\|_{\mathbb{L}^2_{\mu }}= \liminf_{n\to +\infty}\|%
\hat{\rho}(\rho_n)-\bar{\rho}\|_{\mathbb{L}^2_{\mu }}= \hat{\mathcal{P}}%
(\lambda,\bar{\rho}).
\end{equation*}
This indicates that $\hat{\rho}(\lambda_n) \to \breve{\rho}$ in $\mathbb{L}%
^2_{\mu}$ as $n \to +\infty$. The remainder of the proof follows as before.
\end{proof}

\end{document}
